

\maketitle

\begingroup\small
\noindent\textbf{Note:} A companion manuscript develops the divisor-filter mechanism into a general framework, the \href{https://arxiv.org/abs/2509.12297}{Fej\'er--Dirichlet Lift}, and establishes its connection to Dirichlet series and L-functions.
\par\medskip
\endgroup

\begin{abstract}
A $C^1$ prime indicator $\mathcal{P}\colon\mathbb{R}\to\mathbb{R}$ is constructed by applying the Fejér identity to the sine–quotient encoder of trial division. For integers $n\ge 2$, $\mathcal P(n)=0$ holds exactly for odd primes; $\mathcal P(2)>0$. For all non-integers $x>1$ one has $\mathcal P(x)>0$. The function is piecewise $C^\infty$ and its second derivative has jumps precisely at the squares $m^2$, with explicit sizes, all larger than $2$. Replacing the sharp cut-off by a smooth transition yields $C^\infty$ analogues $\mathcal{P}_\tau$ and $\mathcal{P}_\sigma$ with integer limits $\mathcal{P}_\tau(n;\kappa)\to \tau(n)-2$ and $\mathcal{P}_\sigma(n;\kappa)\to \sigma(n)-n-1$ as $\kappa\to\infty$, obtained from locally uniform convergence of derivative series. For large $\kappa$, numerical evidence indicates companion zeros near odd primes for $\mathcal{P}_\tau$ and an asymmetric pair for $\mathcal{P}_\sigma$; for $\mathcal{P}_\sigma$, a computer-assisted argument shows that, for every prime $p\ge5$ and $\kappa\ge 3.25\,(p+1)\ln p$, there is exactly one zero in $(p,p+1)$. The indicator interpretation is confined to integer input, and no statements are claimed about the prime number theorem or zero distributions of $L$-functions. The appendix includes two illustrative prime-counting sums.
\end{abstract}

\clearpage
\begingroup
\hypersetup{linkcolor=black} % local override for ToC
\tableofcontents
\endgroup
% restore global link color for the rest of the paper
\hypersetup{linkcolor=blue}
\bigskip
\clearpage

% ---------------------------------------------------------------------------
